%% Annexe B : Référentiel des outils
\chapter{Référentiel des Outils}
\label{ann:outils}

Les outils et modèles cités sont des exemples. Le cours insiste sur la
méthode~: l'outil sert la décision, il ne la remplace pas.

\begin{tableaucours}{L{4.2cm}L{5.4cm}L{2.8cm}}{%
    \tableauHeader{Outil ou modèle} &
    \tableauHeader{Usage} &
    \tableauHeader{Chapitre, lab}
}
Inventaire des actifs & Recenser et classer l'information & Ch.~\ref{chap:cia}, SGSI-LAB-01 \\
Modèle de politique & Rédiger la PSSI & Ch.~\ref{chap:politique}, SGSI-LAB-02 \\
Grille d'écart ISO 27001 & Mesurer la maturité (gap analysis) & Ch.~\ref{chap:clauses46}, SGSI-LAB-03 \\
Méthode EBIOS RM & Apprécier les risques & Ch.~\ref{chap:methodes}, SGSI-LAB-04 \\
Matrice des risques & Hiérarchiser et traiter & Ch.~\ref{chap:traitement}, SGSI-LAB-05 \\
Modèle de DdA & Relier risques et mesures & Ch.~\ref{chap:dda}, SGSI-LAB-06 \\
Matrice RBAC & Contrôler les accès par rôle & Ch.~\ref{chap:technique}, SGSI-LAB-07 \\
Plan de continuité & Sauvegarde 3-2-1, RTO/RPO & Ch.~\ref{chap:continuite}, SGSI-LAB-08 \\
Tableau de bord & Suivre les indicateurs & Ch.~\ref{chap:indicateurs}, SGSI-LAB-09 \\
Fiche de non-conformité & Conduire un audit interne & Ch.~\ref{chap:audit}, SGSI-LAB-10 \\
\end{tableaucours}

\begin{boxalert}
\textbf{Rappel.} Une mesure ne se retient que si un risque la justifie. Le
catalogue de l'Annexe~A n'est pas une liste à appliquer en bloc~: c'est une
réponse au risque apprécié.
\end{boxalert}
